\documentclass[11pt,a4paper]{article}
\usepackage{auditstyle}
\lhead{Fresh Glimm--Jaffe proof audit}
\rhead{16 August 2026}
\title{Corrected structured proof II\\[3pt]
Uniform gap and local moment $\Longrightarrow$ local-norm convergence and the main theorem}
\author{Fresh reconstruction of \texttt{part\_i/main.tex}}
\date{16 August 2026}

\begin{document}
\maketitle

\begin{resultbox}
\textbf{Status.} \Conditional.  Given the explicitly listed inputs in Section
\ref{sec:inputs}, the full cutoff net is Cauchy in the norm of every local dual,
with a super-polynomial rate.  The resulting state is locally normal, translation
invariant, parity invariant for even interaction, and a spectral ground state.

The derivation below is independent of Euclidean uniqueness.  The current Part I
document correctly states its two-phase exclusion only conditionally on the
load-bearing phase-decomposition hypothesis B4; the final section records that scope.
\end{resultbox}

\tableofcontents

\section{Exact inputs}\label{sec:inputs}

For a spatial cutoff $g\in\G$, let
\[
 K(g)=H_0+V(g),\quad E(g)=\inf\spec K(g),\quad
 H(g)=K(g)-E(g),\quad P_g=|\Om_g\rangle\langle\Om_g|,
\]
and let $\omega_g(A)=\ip{\Om_g}{A\Om_g}$ and
$\beta_t^g=\operatorname{Ad}e^{itH(g)}$.  Let $\A(O)$ be the time-zero local
von Neumann algebra of a bounded interval $O$ and let $O_r$ be its open
$r$-neighbourhood.

\begin{assumption}[FKN, stability, and free hypercontractivity (F)]\label{ass:F}
For every affine cutoff path $g_s=(1-s)g+sg'$ and every $T>0$, the interacting
FKN formula (3.1) below holds.  Its finite-slab interaction random variables satisfy
$\mathcal U_v\in L^p$ for every finite $p$ and
$e^{-a\mathcal U_g}\in L^1$ for every finite $a>0$.  The free time-zero Markov
semigroup satisfies Nelson hypercontractivity with
$(q-1)/(p-1)=e^{2m_0T}$.  These are the precise standard inputs used to prove path
regularity; they are not consequences of (G) or (M).
\end{assumption}

\begin{assumption}[uniform physical gap (G)]\label{ass:G}
There is $\gamma>0$, independent of $g$, such that
\[
 \spec H(g)\subseteq\{0\}\cup[\gamma,\infty),
 \qquad \dim\ker H(g)=1.                                           \tag{G}
\]
\end{assumption}

\begin{assumption}[uniform local interaction moment (M)]\label{ass:M}
There is $M<\infty$ such that
\[
 \|V(h)\Om_g\|\leq M                                             \tag{M}
\]
for every $g$ and every real $h\in C_c^\infty(\R)$ satisfying
$\|h\|_\infty\leq1$ and whose support is contained in an interval of length $2$.
\end{assumption}

\begin{assumption}[local domains and affiliation]\label{ass:D}
For every compactly supported real $h$, $\Om_g\in D(V(h))$, and $V(h)$ is
self-adjoint and affiliated with $\A(U)$ whenever $U$ contains $\supp h$.
Along every affine path there is a common core $\mathscr C$ for the $K_s$ such that
$K_s\xi=K_{s_0}\xi+(s-s_0)V(g'-g)\xi$ for $\xi\in\mathscr C$.
The correct model-specific smoothing source is Glimm--Jaffe Expositions,
Theorem 7.4.  Theorem 7.3 alone assumes that the multiplication potential is
bounded below and therefore does not, by itself, justify this application.
\end{assumption}

\begin{assumption}[finite propagation and patching]\label{ass:P}
For every cutoff,
\[
 \beta_t^g(\A(O))\subseteq\A(O_{|t|}).                             \tag{FP}
\]
There is a cutoff-independent quasi-local dynamics $\alpha$ such that
\[
 g=1\text{ on }O_{|t|}\quad\Longrightarrow\quad
 \alpha_t(A)=\beta_t^g(A)\quad(A\in\A(O)).                         \tag{P}
\]
\end{assumption}

\begin{assumption}[covariance]\label{ass:T}
If $(\tau_ag)(x)=g(x-a)$, then
\[
 \omega_{\tau_ag}=\omega_g\circ\tau_{-a}.                         \tag{T}
\]
For even $\mathscr P$, each $\omega_g$ is invariant under
$\theta=\operatorname{Ad}\Gamma(-1)$.
\end{assumption}

The first corrected proof document derives Assumptions \ref{ass:G} and \ref{ass:M}
from GJS in the weak-coupling region.  Assumptions \ref{ass:F} and
\ref{ass:D}--\ref{ass:T} are standard but logically independent
cutoff-construction inputs.

\section{Statement to be proved}

For a cutoff $g$ and interval $O$, write
\[
 d(g,O)=\dist\bigl(O,\R\setminus\{g=1\}^{\circ}\bigr).
\]

\begin{theorem}[local-norm convergence and spectral ground state]\label{thm:main}
Under Assumptions \ref{ass:F}--\ref{ass:T}, there is a state $\omega_\infty$ on
the quasi-local algebra $\A$ such that, for every bounded interval $O$,
\[
 \lim_{g\in\G}\| (\omega_g-\omega_\infty)|_{\A(O)}\|=0.           \tag{2.1}
\]
If $g=1$ on $O_d$, then
\[
 \|(\omega_g-\omega_\infty)|_{\A(O)}\|\leq\Psi_\gamma(d),         \tag{2.2}
\]
where $\Psi_\gamma(d)=O_N((1+d)^{-N})$ for every $N$.  The limit is locally
normal, is a spectral ground state for $(\A,\alpha)$, is translation invariant,
and, for even $\mathscr P$, is $\theta$-invariant.
\end{theorem}

\section{Structured proof}

\LS{1}{1}{Prove differentiability of the ground-state projection along an affine
cutoff path.}

Let $g,g'$ equal $1$ on a common nonempty interval, put
\[
 g_s=(1-s)g+sg',\quad v=g'-g,\quad K_s=K(g_s),
 \qquad H_s=K_s-E_s,\\
 P_s=|\Om_s\rangle\langle\Om_s|,\qquad Q_s=1-P_s.
\]

\LS{2}{1}{The family $s\mapsto e^{-TK_s}$ is norm-continuous on $[0,1]$ and
norm-$C^1$ on $(0,1)$ for each $T>0$.}

\Because On free Euclidean path space define
\[
 \mathcal U_s=\lambda\int_0^T\!\int g_s(x){:}\mathscr P(\Phi(t,x)){:}\,dxdt,
 \quad w_s=e^{-\mathcal U_s},\quad \mathcal U_v=\mathcal U_1-\mathcal U_0.
\]
For bounded time-zero $f,h$, the interacting FKN formula is
\[
 \ip{f}{e^{-TK_s}h}=\int\overline{J_Tf}\,J_0h\,w_s\,d\mu_0.        \tag{3.1}
\]
The correct references are Simon Corollary V.14/Theorem V.15 or GRS Theorem
II.16, not Simon Theorem III.12.

Set $a=e^{-m_0T}$ and
\[
 q_T=\frac2{1-a},\quad r_T=\frac2{1+a},\quad p_T=1+a,
 \quad p_T'=1+a^{-1}.
\]
Then $1/r_T+1/q_T=1$, $r_Tp_T=2$, and
$(p_T'-1)/(p_T-1)=e^{2m_0T}$.  Nelson hypercontractivity---Simon
Theorem I.17, or its time-region version Theorem III.17---therefore gives
\[
 \|\mathsf P_TF\|_{p_T'}\leq\|F\|_{p_T}.
\]
The Markov property and H\"older yield
\begin{align*}
 \|(J_Tf)(J_0h)\|_{r_T}^{r_T}
 &=\int |h|^{r_T}\mathsf P_T(|f|^{r_T})d\nu_0\\
 &\leq\||h|^{r_T}\|_{p_T}\|\mathsf P_T(|f|^{r_T})\|_{p_T'}\\
 &\leq\|h\|_2^{r_T}\|f\|_2^{r_T}.
\end{align*}
Thus
\[
 \|(J_Tf)(J_0h)\|_{r_T}\leq\|f\|_2\|h\|_2.                       \tag{3.2}
\]

Finite-volume stability gives $\mathcal U_v\in L^p(\mu_0)$ for every finite
$p$ and $e^{-a\mathcal U_i}\in L^1(\mu_0)$ for every $a<\infty$, $i=0,1$
(Simon Theorem V.7).  Since
$\mathcal U_s=(1-s)\mathcal U_0+s\mathcal U_1$, H\"older gives
\[
 \|w_s\|_a^a\leq
 \left(\int e^{-a\mathcal U_0}d\mu_0\right)^{1-s}
 \left(\int e^{-a\mathcal U_1}d\mu_0\right)^s.                    \tag{3.3}
\]
Convergence in probability together with (3.3) at a larger exponent gives
$w_{s_n}\to w_s$ in $L^a$ by Vitali.  Moreover
\[
 \frac{w_{s+h}-w_s}{h}
 =-\mathcal U_v\int_0^1w_{s+\vartheta h}\,d\vartheta.              \tag{3.4}
\]
Choose finite $p,a$ with $q_T^{-1}=p^{-1}+a^{-1}$.  H\"older and the just-proved
$L^a$ continuity show that (3.4) converges in $L^{q_T}$ to
$-\mathcal U_vw_s$ for $0<s<1$.  Finally (3.1), (3.2), and duality give
\[
 \|e^{-TK_s}-e^{-TK_r}\|\leq\|w_s-w_r\|_{q_T},
\]
and the same estimate for difference quotients.  This proves the asserted norm
regularity.

\LS{2}{2}{$s\mapsto P_s$ is norm-continuous on $[0,1]$ and $C^1$ on $(0,1)$.}

\Because $L_s=e^{-TK_s}$ has the simple largest eigenvalue
$\lambda_s=e^{-TE_s}=\|L_s\|$.  By (G),
\[
 \spec L_s\setminus\{\lambda_s\}\subset[0,e^{-T\gamma}\lambda_s].
\]
The relative spectral separation $1-e^{-T\gamma}$ is uniform.  A Riesz contour
around $\lambda_s$, together with the norm regularity from $\langle2\rangle1$,
gives exactly the stated regularity of $P_s$.

\LS{2}{3}{Locally on $(0,1)$ choose normalized $C^1$ vectors in the parallel gauge
$\ip{\Om_s}{\dot\Om_s}=0$.  Then}
\[
 \dot E_s=\ip{\Om_s}{V(v)\Om_s},\qquad
 \dot\Om_s=-H_s^{-1}Q_sV(v)\Om_s.                                 \tag{3.5}
\]

\Because Let $\mathscr C$ be the common core supplied by cutoff construction.
For $\xi\in\mathscr C$,
$K_s\xi=K_{s_0}\xi+(s-s_0)V(v)\xi$.  Differentiate the weak eigenvalue identity
$\ip{K_s\xi}{\Om_s}=E_s\ip{\xi}{\Om_s}$ to obtain
\[
 \ip{H_s\xi}{\dot\Om_s}=\ip{\xi}{(\dot E_s-V(v))\Om_s}.           \tag{3.6}
\]
Assumption \ref{ass:D} makes the vector on the right well defined.  Since
$\mathscr C$ is a core for the self-adjoint $H_s$, (3.6) says
$\dot\Om_s\in D(H_s)$ and
$H_s\dot\Om_s=(\dot E_s-V(v))\Om_s$.  Pairing with $\Om_s$ proves the first
formula in (3.5).  The parallel gauge puts $\dot\Om_s$ in $Q_s\Hh$; (G) makes
$H_s^{-1}Q_s$ bounded by $\gamma^{-1}$, proving the second formula.

\LS{2}{4}{This proves $\langle1\rangle1$.}\QedStep

\LS{1}{2}{Construct an exact inverse-energy filter with super-polynomial time tail.}

\LS{2}{1}{Choose even $\chi\in C^\infty(\R;[0,1])$ with
$\chi=0$ on $[-\gamma/2,\gamma/2]$ and $\chi=1$ outside
$(-\gamma,\gamma)$, and set}
\[
 \widehat W_\gamma(E)=\begin{cases}\chi(E)/E,&E\ne0,\\0,&E=0.\end{cases} \tag{3.7}
\]

\LS{2}{2}{Its inverse transform
$W_\gamma(t)=(2\pi)^{-1}\int e^{-itE}\widehat W_\gamma(E)dE$ is an
$L^1$ function satisfying}
\[
 \int(1+|t|)^q|W_\gamma(t)|dt<\infty\quad(q\geq0).                 \tag{3.8}
\]

\Because Oddness gives
\[
 W_\gamma(t)=-\frac{i}{\pi}\int_0^\infty\frac{\sin(tE)\chi(E)}E\,dE.
\]
For $0<|t|\leq1$, the compact transition part is absolutely bounded and the tail is
uniformly bounded after $u=|t|E$ by Dirichlet's test.  For $|t|\geq1$, every
derivative $\widehat W_\gamma^{(k)}$ with $k\geq1$ is integrable, so $k$
integrations by parts give
\[
 |W_\gamma(t)|\leq\frac{\|\widehat W_\gamma^{(k)}\|_1}{2\pi|t|^k}.
\]
Choose $k>q+1$.  This proves (3.8), including integrability near zero.

\LS{2}{3}{If $H$ satisfies (G), then}
\[
 H^{-1}(1-P_{\ker H})=\int_\R W_\gamma(t)e^{itH}(1-P_{\ker H})\,dt. \tag{3.9}
\]

\Because Fourier inversion gives $\widehat W_\gamma(E)=E^{-1}$ on
$[\gamma,\infty)$ and zero at $E=0$.  The spectral theorem gives (3.9); it is a
norm-convergent Bochner integral because $W_\gamma\in L^1$.

Define
\[
 T_\gamma(r)=\int_{|t|>r}|W_\gamma(t)|dt.
\]
From (3.8), Markov's inequality gives
$T_\gamma(r)\leq C_{q,\gamma}(1+r)^{-q}$ for every $q$.

\LS{2}{4}{This proves $\langle1\rangle2$.}\QedStep

\LS{1}{3}{Prove the spectral-separation estimate used for clustering.}

Suppose
\[
 F(t)=\int_{[\gamma,\infty)}e^{itE}d\rho(E),\qquad
 G(t)=\int_{(-\infty,-\gamma]}e^{itE}d\sigma(E),
\]
with $|\rho|(\R),|\sigma|(\R)\leq M_0$, and $F=G$ on $(-R,R)$, where
$\gamma R\geq1$.

\LS{2}{1}{For $a,\varepsilon>0$ define}
\[
 f_{a,\varepsilon}(t)=\frac{1}{2\pi i}
 \frac{e^{-t^2/(2a)}}{t-i\varepsilon}.
\]
Its transform is
\[
 \check f_{a,\varepsilon}(E)=
 \int_0^\infty\sqrt{\frac a{2\pi}}
 e^{-a(E-E')^2/2}e^{-\varepsilon E'}dE',                           \tag{3.10}
\]
which lies in $[0,1]$ and tends to the standard normal distribution function
$\Phi(\sqrt aE)$ as $\varepsilon\downarrow0$.

\Because The transform of $(2\pi i)^{-1}(t-i\varepsilon)^{-1}$ is
$\mathbf1_{(0,\infty)}(E)e^{-\varepsilon E}$ by a one-pole contour integral.
Multiplication by the Gaussian in $t$ convolves this function with the normalized
Gaussian displayed in (3.10).

\LS{2}{2}{The part of $f_{a,\varepsilon}$ outside $(-R,R)$ satisfies}
\[
 \int_{|t|\geq R}|f_{a,\varepsilon}(t)|dt
 \leq\frac a{\pi R^2}e^{-R^2/(2a)}.                               \tag{3.11}
\]

\Because $|t-i\varepsilon|\geq|t|$ and, for $t\geq R$,
$t^{-1}\leq t/R^2$.  Integrating $te^{-t^2/(2a)}$ exactly proves (3.11).

\LS{2}{3}{Using $F=G$ on $(-R,R)$ and Fubini gives}
\[
 \left|\int\Phi(\sqrt aE)d\rho(E)-
       \int\Phi(\sqrt aE)d\sigma(E)\right|
 \leq\frac{2M_0a}{\pi R^2}e^{-R^2/(2a)}.                          \tag{3.12}
\]

\Because Before $\varepsilon\downarrow0$, the difference of the two time integrals
is supported in $|t|\geq R$ and is at most
$(|\rho|(\R)+|\sigma|(\R))$ times (3.11).  Dominated convergence in energy proves
(3.12).

\LS{2}{4}{Therefore}
\[
 |F(0)|\leq2M_0e^{-\gamma R/2}.                                   \tag{3.13}
\]

\Because Insert and subtract the two integrals in (3.12).  On the supports,
$1-\Phi(\sqrt aE)\leq\tfrac12e^{-a\gamma^2/2}$ and
$\Phi(\sqrt aE)\leq\tfrac12e^{-a\gamma^2/2}$, respectively.  Hence
\[
 |F(0)|\leq M_0e^{-a\gamma^2/2}
 +\frac{2M_0a}{\pi R^2}e^{-R^2/(2a)}.
\]
Set $a=R/\gamma$.  Both exponentials become $e^{-\gamma R/2}$ and
$2a/(\pi R^2)=2/(\pi\gamma R)\leq2/\pi<1$, which proves (3.13).

\LS{2}{5}{This proves $\langle1\rangle3$.}\QedStep

\LS{1}{4}{Derive clustering with one unbounded local observable.}

Let $H,\Om$ satisfy (G), let $X=X^*$ be affiliated with every $\A(U)$ containing
a compact set $B$, assume $\Om\in D(X)$, and let $A\in\A(O)$.  Put
$R=\dist(B,O)>0$ and suppose $\gamma R\geq1$.

\LS{2}{1}{Define}
\begin{align*}
 F(t)&=\ip{X\Om}{(e^{itH}-P_0)A\Om},\\
 G(t)&=\ip{A^*\Om}{(e^{-itH}-P_0)X\Om},
 \qquad P_0=|\Om\rangle\langle\Om|.
\end{align*}
Then $F$ and $G$ have spectral measures supported in
$[\gamma,\infty)$ and $(-\infty,-\gamma]$, with variations at most
$M_0=\|X\Om\|\|A\|$.

\Because Insert the spectral resolution of $H$ on $(1-P_0)\Hh$ and use
Cauchy--Schwarz for the variation of the resulting complex measures.

\LS{2}{2}{$F(t)=G(t)$ for $|t|<R$.}

\Because By finite propagation, $\beta_t(A)\in\A(O_{|t|})$.  Choose a bounded
open $U\supset B$ disjoint from $O_{|t|}$.  Affiliation implies that
$\beta_t(A)\in\A(U)'$ preserves $D(X)$ and commutes with $X$ there.  Since
$e^{-itH}\Om=\Om$,
\begin{align*}
 \ip{X\Om}{e^{itH}A\Om}
 &=\ip{\Om}{X\beta_t(A)\Om}
  =\ip{\Om}{\beta_t(A)X\Om}\\
 &=\ip{A^*\Om}{e^{-itH}X\Om}.
\end{align*}
The $P_0$ terms agree because $\ip{\Om}{X\Om}$ is real.

\LS{2}{3}{Apply (3.13) to obtain}
\[
 \left|\ip{X\Om}{A\Om}-\ip{X\Om}{\Om}\ip{\Om}{A\Om}\right|
 \leq2\|X\Om\|\|A\|e^{-\gamma R/2}.                              \tag{3.14}
\]

\LS{2}{4}{This proves $\langle1\rangle4$.}\QedStep

\LS{1}{5}{Differentiate the state and localize the perturbation.}

\LS{2}{1}{For self-adjoint $A\in B(\Hh)$ and $0<s<1$,}
\[
 \frac d{ds}\omega_{g_s}(A)
 =-2\operatorname{Re}\int_\R W_\gamma(t)
 \left[\ip{V(v)\Om_s}{\beta_t^{g_s}(A)\Om_s}
 -\ip{V(v)\Om_s}{\Om_s}\omega_{g_s}(A)\right]dt.                 \tag{3.15}
\]

\Because Differentiate the vector state and insert (3.5):
\[
 \frac d{ds}\omega_{g_s}(A)
 =-2\operatorname{Re}\ip{V(v)\Om_s}{H_s^{-1}Q_sA\Om_s}.
\]
Insert the exact filter (3.9).  Since
$e^{itH_s}Q_sA\Om_s=\beta_t^{g_s}(A)\Om_s-omega_{g_s}(A)\Om_s$,
this is (3.15).  The integral is absolutely convergent because
$W_\gamma\in L^1$ and $V(v)\Om_s$ exists; $v$ has compact support.

\LS{2}{2}{Choose a smooth partition of unity $\sum_j\eta_j=1$ with
$0\leq\eta_j\leq1$ and $\supp\eta_j\subset(j-1,j+1)$.  Put
$h_j=\eta_jv$, $V_j=V(h_j)$.  Then}
\[
 V(v)\xi=\sum_jV_j\xi,\qquad \|V_j\Om_s\|\leq M.                  \tag{3.16}
\]

\Because Only finitely many $h_j$ are nonzero.  Also $|v|\leq1$, so
$\|h_j\|_\infty\leq1$, and each support has length at most $2$.  Apply (M) to
the cutoff $g_s$.

\LS{2}{3}{This proves $\langle1\rangle5$.}\QedStep

\LS{1}{6}{Prove the one-path local-perturbations estimate.}

Assume $g=g'=1$ on $O_d$, where
\[
 d\geq d_0:=4+2/\gamma.
\]
For a nonzero $h_j$ set $B_j=\supp h_j$ and $r_j=\dist(B_j,O)$.

\LS{2}{1}{$r_j\geq d$ for every nonzero $h_j$.}

\Because On $O_d$, $v=g'-g=0$.  Therefore the support of $h_j=\eta_jv$ is
contained in $\R\setminus O_d$.

\LS{2}{2}{For self-adjoint $A\in\A(O)$, $\|A\|\leq1$, define}
\[
 C_{j,s}(t)=\ip{V_j\Om_s}{\beta_t^{g_s}(A)\Om_s}
 -\ip{V_j\Om_s}{\Om_s}\omega_{g_s}(A).
\]
Then
\[
 |C_{j,s}(t)|\leq
 \begin{cases}
 2Me^{-\gamma r_j/4},&|t|\leq r_j/2,\\
 2M,&t\in\R.
 \end{cases}                                                       \tag{3.17}
\]

\Because In the first case, finite propagation places
$\beta_t^{g_s}(A)$ in $\A(O_{|t|})$ and
$\dist(B_j,O_{|t|})\geq r_j/2$.  Moreover
$\gamma r_j/2\geq1$.  Apply (3.14) with $X=V_j$, then (3.16).
For arbitrary $t$, bound each of the two terms defining $C_{j,s}(t)$ by $M$.

\LS{2}{3}{Equations (3.15)--(3.17) imply}
\[
 \left|\frac d{ds}\omega_{g_s}(A)\right|
 \leq4M\sum_j\left(\|W_\gamma\|_1e^{-\gamma r_j/4}
                    +T_\gamma(r_j/2)\right).                      \tag{3.18}
\]

\Because Split the $t$-integral for each $j$ at $|t|=r_j/2$.  The outside
factor $2$ in (3.15), the factor $2M$ in either line of (3.17), and the two
integral regions give exactly the coefficient $4M$ in (3.18).

\LS{2}{4}{At most eight indices $j$ have $r_j\in[n,n+1)$ for any integer
$n\geq\lfloor d\rfloor$.}

\Because Write $O=(a,b)$.  Since $d\geq4$ and
$B_j\subset[j-1,j+1]$, a nonempty $B_j$ lies wholly on one side of $O_d$.
On the right, compactness gives $x_j\in B_j$ with $x_j-b=r_j$.  Since
$|x_j-j|\leq1$ and $r_j\in[n,n+1)$,
\[
 j\in[b+n-1,b+n+2],
\]
which contains at most four integers.  The left side gives four more.

\LS{2}{5}{Integrating (3.18) along the path gives}
\[
 \|(\omega_g-\omega_{g'})|_{\A(O)}\|
 \leq32M\sum_{n\geq\lfloor d\rfloor}
 \left(\|W_\gamma\|_1e^{-\gamma n/4}+T_\gamma(n/2)\right).        \tag{3.19}
\]

\Because Both functions of $r_j$ in (3.18) decrease with $r_j$.  Apply the shell
count from $\langle2\rangle4$, producing $4M\cdot8=32M$.  Integrate first on
$[\varepsilon,1-\varepsilon]$, where differentiability holds, and then let
$\varepsilon\downarrow0$ using the endpoint norm continuity of $P_s$ proved in
$\langle1\rangle1$.  Taking the supremum over self-adjoint contractions gives the
norm of the hermitian functional $\omega_g-\omega_{g'}$.

\LS{2}{6}{This proves $\langle1\rangle6$.}\QedStep

\LS{1}{7}{Define the rate and prove that it is super-polynomial.}

For $d\geq d_0$ set
\[
 \Psi_\gamma(d)=32M\sum_{n\geq\lfloor d\rfloor}
 \left(\|W_\gamma\|_1e^{-\gamma n/4}+T_\gamma(n/2)\right),         \tag{3.20}
\]
and for $d<d_0$ enlarge it to at least $2$.

\LS{2}{1}{For every $N\geq1$, there is $C_{N,\gamma,M}$ such that}
\[
 \Psi_\gamma(d)\leq C_{N,\gamma,M}(1+d)^{-N}.                     \tag{3.21}
\]

\Because The first series has an exponential tail.  For the second choose
$q=N+2$ in $T_\gamma(r)\leq C_{q,\gamma}(1+r)^{-q}$; then
$\sum_{n\geq\lfloor d\rfloor}(1+n)^{-N-2}=O((1+d)^{-N-1})$.

\LS{2}{2}{This proves $\langle1\rangle7$.}\QedStep

\LS{1}{8}{Prove the full cutoff net is locally norm Cauchy and construct
$\omega_\infty$.}

\LS{2}{1}{Fix $O$ and $\varepsilon>0$.  Choose $d\geq d_0$ with
$\Psi_\gamma(d)<\varepsilon$, and choose $g_0\in\G$ with $g_0=1$ on $O_d$.}

\LS{2}{2}{If $g,g'\succcurlyeq g_0$, then}
\[
 \|(\omega_g-\omega_{g'})|_{\A(O)}\|<\varepsilon.                 \tag{3.22}
\]

\Because The directed order is inclusion of plateau interiors, so both cutoffs equal
$1$ on $O_d$.  Apply (3.19).

\LS{2}{3}{The limit functionals on local algebras are compatible and extend uniquely
to a state $\omega_\infty$ on the norm-dense quasi-local union.}

\Because Each $\A(O)^*$ is complete, so (3.22) gives a limit on $\A(O)$.  If
$O_1\subset O_2$, uniqueness of the norm limit shows that the $O_2$ limit restricts
to the $O_1$ limit.  Positivity and normalization pass to the limit.  A compatible
bounded functional of norm one on the algebraic union extends uniquely by norm
continuity to $\A$.

\LS{2}{4}{If one fixed $g$ equals $1$ on $O_d$, then (2.2) holds.}

\Because The subset of cutoffs $g'$ equal to $1$ on $O_d$ is cofinal.  Apply
(3.19) to $(g,g')$ and let $g'$ tend through this subset to the local norm limit.

\LS{2}{5}{This proves local-norm convergence and its rate.}\QedStep

\LS{1}{9}{Prove local normality.}

Fix $O$ and let $0\leq A_i\uparrow A$ in $\A(O)$.  For any cutoff $g$,
\[
 0\leq\omega_\infty(A)-\omega_\infty(A_i)
 \leq2\delta\|A\|+\omega_g(A-A_i)                                 \tag{3.23}
\]
whenever $\|(\omega_g-\omega_\infty)|_{\A(O)}\|<\delta$.

\Because Add and subtract $\omega_g(A)$ and $\omega_g(A_i)$, use
$\|A_i\|\leq\|A\|$, and bound each state difference by the local dual norm.
The cutoff vector state is normal on $\A(O)$, so its last term in (3.23) tends to
zero with $i$.  First let $i$ increase, then let $\delta\downarrow0$.  Thus
$\omega_\infty$ is normal on $\A(O)$.

\LS{1}{10}{Prove the spectral ground-state condition.}

Choose cutoffs $g_n$ with $g_n=1$ on $(-n,n)$.  For local
$A,B\in\A(O)$ define the finite complex measure
\[
 \rho_n(S)=\ip{A^*\Om_{g_n}}{\mathbf1_S(H(g_n))B\Om_{g_n}},
 \quad G_n(t)=\int_{[0,\infty)}e^{itE}d\rho_n(E).                  \tag{3.24}
\]

\LS{2}{1}{For $|t|<d(g_n,O)$,}
\[
 \omega_{g_n}(A\alpha_t(B))=G_n(t).                               \tag{3.25}
\]

\Because Patching gives $\alpha_t(B)=\beta_t^{g_n}(B)$, and
$e^{-itH(g_n)}\Om_{g_n}=\Om_{g_n}$.  Moving the unitary to the middle gives (3.25).

\LS{2}{2}{If $f\in L^1(\R)$ and
$\check f(E)=\int f(t)e^{itE}dt$ vanishes on $[0,\infty)$, then}
\[
 \int_\R f(t)G_n(t)dt=0.                                         \tag{3.26}
\]

\Because Fubini is justified by
$|\rho_n|(\R)\leq\|A\|\|B\|$, and the spectral support in (3.24) is nonnegative.

\LS{2}{3}{The limit function $G(t)=\omega_\infty(A\alpha_t(B))$ is continuous and
satisfies $\int fG=0$.}

\Because For each $T$, eventually $d(g_n,O)>T$, and
$A\alpha_t(B)\in\A(O_T)$ for $|t|\leq T$.  Hence
\[
 \sup_{|t|\leq T}|G(t)-G_n(t)|
 \leq\|A\|\|B\|\,
 \|(\omega_\infty-\omega_{g_n})|_{\A(O_T)}\|\longrightarrow0.     \tag{3.27}
\]
Thus $G$ is a locally uniform limit of continuous functions.  Let
$M_0=\|A\|\|B\|$ and $\epsilon_T=\int_{|t|>T}|f(t)|dt$.  Using (3.26),
\begin{align*}
 \left|\int fG\right|
 &\leq \|f\|_1M_0
 \|(\omega_\infty-\omega_{g_n})|_{\A(O_T)}\|+2M_0\epsilon_T.
\end{align*}
First let $n\to\infty$, then $T\to\infty$.

\LS{2}{4}{$\omega_\infty\circ\alpha_t=\omega_\infty$.}

\Because For local $A$ and fixed $t$, eventually patching gives
$\omega_{g_n}(\alpha_t(A))=\omega_{g_n}(A)$.  Both observables lie in a fixed
local algebra, so local norm convergence passes the identity to the limit.  Density
extends it to $\A$.

\LS{2}{5}{Steps $\langle2\rangle3$--$\langle2\rangle4$ are exactly the spectral
ground-state condition.}\QedStep

\LS{1}{11}{Prove spatial and parity invariance.}

\LS{2}{1}{$\omega_\infty\circ\tau_a=\omega_\infty$ for every $a\in\R$.}

\Because $g\mapsto\tau_ag$ is an order isomorphism of the cutoff net and hence is
cofinal.  For local $A$, covariance (T) and local norm convergence give
\[
 \omega_\infty(A)=\lim_g\omega_{\tau_ag}(A)
 =\lim_g\omega_g(\tau_{-a}A)=\omega_\infty(\tau_{-a}A).
\]
Extend by norm density.

\LS{2}{2}{For even $\mathscr P$, $\omega_\infty\circ\theta=\omega_\infty$.}

\Because Every cutoff state has this invariance by Assumption \ref{ass:T}; pass to
the local norm limit, then use density.

\LS{2}{3}{Together with $\langle1\rangle8$--$\langle1\rangle10$, this proves
Theorem \ref{thm:main}.}\QedStep

\section{The averaged Glimm--Jaffe vacuum}

The manuscript's canonical-vacuum corollary does follow.  Let
$g_n(x)=g(x/n)$ with $g=1$ on $[-3,3]$, and let
$h\geq0$, $\supp h\subset[-1,1]$, $\int h=1$.  The averaged state of Glimm--Jaffe
III has the form
\[
 \omega_n(A)=\int_{-1}^1h(u)\omega_{\tau_{-\epsilon nu}g_n}(A)du,
 \qquad \epsilon\in\{\pm1\}.
\]
For $O\subset[-R,R]$, every shifted cutoff in the integral equals $1$ on an
interval containing $[-2n,2n]$, so
\[
 d(\tau_{-\epsilon nu}g_n,O)\geq2n-R
\]
uniformly in $u$.  Therefore
\[
 \|(\omega_n-\omega_\infty)|_{\A(O)}\|
 \leq\int_{-1}^1h(u)\Psi_\gamma(2n-R)du
 =\Psi_\gamma(2n-R)\longrightarrow0.
\]
No subnet and no pre-established translation invariance of the limit is used.

\section{The conditional two-phase exclusion}

The current Part I theorem correctly makes the following four assumptions explicit:

\begin{enumerate}[label=\textup{(B\arabic*)}]
\item there exist two translation-invariant clustering ground states
$\omega_+,\omega_-$;
\item $\theta\omega_+=\omega_-$;
\item a local order parameter $A=A^*$ has $\omega_\pm(A)=\pm m$, $m\ne0$;
\item every $\theta$-invariant ground state relevant there is the midpoint
$\tfrac12(\omega_++\omega_-)$.
\end{enumerate}

The contradiction uses (B4) to identify the locally obtained invariant limit with that
midpoint.  Without (B4), invariance alone does not imply the identification.  GJS
1975--1976 prove phase-transition and pure-phase results in specified strong-coupling
models, but the cited pages do not state all four assumptions as a classification of
the Hamiltonian ground-state space.  Hence the exclusion is rigorous exactly as a
conditional theorem and is not an unconditional consequence of Theorem~\ref{thm:main}.

\section{Dependency conclusion}

The exact dependency chain is
\begin{align*}
 \boxed{(F)+(G)+(M)+(D)+(FP)+(P)+(T)}
 &\Longrightarrow
 \boxed{\begin{gathered}\text{path derivative + filter}\\
                         \text{+ unbounded clustering}\end{gathered}}\\
 &\Longrightarrow \boxed{\text{one-path estimate}}\\
 &\Longrightarrow \boxed{\text{full-net local-norm limit}}\\
 &\Longrightarrow
 \boxed{\begin{gathered}\text{normal spectral ground state}\\
                         \text{+ invariances}\end{gathered}}.
\end{align*}
Every arrow has been derived above.  The theorem remains conditional on all seven input groups;
it is not logically ``conditional only'' on the two GJS estimates.

\end{document}
